\section{实战案例与问题诊断}
\subsection{Self-Feeding Chatbot 模板}
Weston 再次回顾 Self-Feeding Chatbot 的实践：模型与用户对话后，把初始 draft 变成 verified response，再把高质量回答补充进训练集，用来支持 self-rewarding（SRT 774-884）。这个环节强调 ``rewarding itself following instructions''，模型会把 verified response 视为另一种 label，并在 future iterations 中予以强化。

\begin{table}[H]
\centering
\begin{tabular}{p{4cm}p{10cm}}
\toprule
阶段 & 描述 \\
\midrule
User Dialogue & 模型与用户互动，记录输入、输出与上下文，形成即刻监督信号。 \\
Model Draft & 模型输出 draft 并生成 Chain-of-Thought，作为判断基础。 \\
Self-Verification & 模型再次审视 draft，生成 verified response，强调逻辑一致性。 \\
Verified Response & 该回答被当作高质量样本供 reward model 学习。 \\
Training Update & 把 verified response 合并到训练集中，持续优化 reward / meta judge。 \\
\bottomrule
\end{tabular}
\caption{Self-Feeding 模板中的关键阶段与数据流。}
\end{table}

\subsection{错误模式与告警}
在 self-rewarding 迭代中常见三类问题：1）模型 overthinking，即越思考越偏离答案；2）hallucination 与自信叠加；3）meta judge drift，比如 judge 内部 calibration drift 导致 verified response 也被误判为差。Weston 也提到 ``we can use judge to cross check itself and see that it was wrong''（SRT 1188-1190），提示我们需要维护一个可信的 check-list 来捕捉异常。

\begin{table}[H]
\centering
\begin{tabular}{p{4cm}p{5cm}p{5cm}}
\toprule
错误模式 & 典型表现 & 应对策略 \\
\midrule
Overthinking & Chain-of-Thought 逐渐变长但答案无提升 & 限制 reasoning depth、引入 meta judge 纠偏 \\
Hallucination + Confidence & 模型自信但回答与事实脱节 & 多元 evaluator、增加自我校验 \\
Meta judge drift & Verified response 也被判为差 & 定期重新校准 SelfT judge 数据 \\
\bottomrule
\end{tabular}
\caption{实战中 Self-Rewarding 需要监控的核心告警。}
\end{table}

\subsection{本章小结}
Self-Feeding Chatbot 提供线上验证模板，错误模式表格则提醒我们：需要通过 meta judge 重新校准、限制 reasoning depth 并增强 hallucination 监控，才能持续保持 self-improvement 在正确方向上。

